\documentclass[11pt]{article}
\usepackage[margin=1in]{geometry}
\usepackage{amsmath,amssymb}
\usepackage{graphicx}
\usepackage{booktabs}
\usepackage{multirow}
\usepackage{hyperref}
\usepackage{enumitem}
\usepackage{xcolor}
\usepackage{todonotes}
\usepackage{caption}
\usepackage{subcaption}
\usepackage{natbib}

% Color scheme and TODO markers
\definecolor{todocolor}{RGB}{200,50,50}
\newcommand{\TODOEXP}[1]{\textcolor{todocolor}{\textbf{[TODO_EXPERIMENT: #1]}}}
\newcommand{\TODOVERIFY}[1]{\textcolor{todocolor}{\textbf{[TODO_VERIFY: #1]}}}
\newcommand{\TODOCHECK}{\textcolor{todocolor}{\textbf{[TODO_CHECK]}}}
\newcommand{\TODO}[1]{\textcolor{todocolor}{\textbf{[TODO: #1]}}}

\title{Reliability of Multi-Source \textit{Spartina alterniflora} Distribution Maps
for Coastal Monitoring in China: Spatial Agreement, Uncertainty, and Implications
for Change Assessment}

\author{Anonymous Reviewer Submission \\
TODO\_AFFILIATION, TODO\_EMAIL}

\date{}

\begin{document}
\maketitle

\begin{abstract}
\textit{Spartina alterniflora} (\textit{Spartina alterniflora}) is one of the most invasive coastal
plants in China, yet national distribution estimates derived from different remote
sensing products have never been systematically harmonized or compared.
This study implements a transparent, reproducible framework to quantify spatial
agreement, area consistency, and sensitivity across three independently produced
Chinese \textit{Spartina alterniflora} distribution products covering overlapping
years (2017--2021). We compare (1) the China Coastal Saltmarsh Atlas (CMSA), a
10\,m annual vector product (2017--2021, UTM Zone 50N); (2) the Institute of
Geographic Sciences (GEODATA) 30\,m raster time series (1990, 2000, 2015, 2020,
Krasovsky\,1940\,Albers equal-area); and (3) the Capital Normal University
CM-SSM sub-meter patch product (2020, 148\,072 individual patches).
We document harmonization protocols, area comparisons under equal-area projections,
cross-product 2020 overlap analysis, and sensitivity to resolution and processing
choices. National 2020 area estimates range from 519.89\,km$^2$ (GEODATA\,30\,m)
to 593.71\,km$^2$ (CM-SSM sub-meter), with CMSA at 571.35\,km$^2$.
Expressing as a percentage of the minimum source-reported area,
$(593.71-519.89)/519.89 \times 100 = 14.2\,\%$, this is the observed
cross-product area range for the same reference year. Values are
source-reported statistics from each product; harmonized rasterized
areas (Section~\ref{sec:results-e3}) are used for IoU/Dice calculation
on the shared EPSG:3410 domain.
We distinguish \textit{inter-product agreement} from \textit{accuracy against
ground truth}; the former quantifies whether independent producers identify
the same locations, the latter remains limited by the absence of nationally
representative gold-standard reference samples. We discuss implications for
historical change assessment, where a single-product CMSA 2017--2021 trend of
+18.6\,\% contrasts with a GEODATA 1990--2020 trend that peaks in 2015 and
declines thereafter.
Pair-wise 2020 spatial agreement (30\,m harmonized raster, EPSG:3410)
shows substantial disagreement even between the two most comparable products:
CMSA vs.\ CM-SSM IoU = 0.5375, Dice = 0.6992. CMSA vs.\ GEODATA (10\,m vs.\
30\,m) drops to IoU = 0.4437, Dice = 0.6147. CM-SSM vs.\ GEODATA shows the
lowest agreement (IoU = 0.3423, Dice = 0.5101). Resolution sensitivity across
30--100\,m is extremely stable ($\leq$0.4\,\% IoU change), indicating that
thematic encoding and patch-delineation differences---not spatial resolution
per se---account for the observed disagreement magnitude. Regional
stratification (RQ2, Section~\ref{sec:results-e5}) reveals 3\,$\times$ IoU
variation across coastal subregions for the CMSA--GEODATA pair.
See Sections~\ref{sec:results-e5} and~\ref{sec:robustness} for regional and
processing-robustness analyses.
\end{abstract}

\textbf{Keywords:} \textit{Spartina alterniflora}; coastal invasive plants;
map reliability; spatial agreement; Intersection over Union; Dice coefficient;
remote sensing uncertainty; multi-source harmonization; coastal monitoring.

\section{Introduction}
\label{sec:introduction}

\subsection{Ecological context}

\textit{Spartina alterniflora} (smooth cordgrass; \textit{Spartina alterniflora}) is a North American
salt marsh grass introduced to China in 1979 for coastal protection and
reclamation\citep{an_1985_spartina}.
Its subsequent uncontrolled expansion across the Yellow Sea, East China Sea,
and South China Sea coastlines has been one of the most extensively documented
invasion events in coastal ecology\citep{li_2009_spartina}. By displacing native
species (notably \textit{Suaeda salsa} and \textit{Scirpus planiculmis}),
altering sediment dynamics, and changing habitat structure, \textit{S.~alterniflora}
has been classified as one of China's top 100 most harmful invasive species
\citep{zhou_2009_invasive}.

National eradication programmes targeting \textit{S.~alterniflora} have been
deployed since the 2010s, and post-eradication recurrence monitoring is now a
policy priority\citep{cma_2021_eradication}. All such programmes depend on
reliable, consistent maps of current and historical \textit{S.~alterniflora}
distribution as their baseline.

\subsection{The multi-source map landscape}

At least four independent Chinese national \textit{S.~alterniflora} distribution
products have been published since 2010, produced by different institutions
using different sensor suites, classification methods, spatial resolutions, and
reference epochs:

\begin{enumerate}[nosep]
\item \textbf{China Coastal Saltmarsh Atlas (CMSA)}: 10\,m annual vector product
  (2017--2021), derived from Sentinel-2 imagery by the Chinese Academy of
  Sciences coastal team\citep{TODO_CMSA_CITATION}.
\item \textbf{Institute of Geographic Sciences (IGSNRR) GEODATA time series}:
  30\,m raster layers for 1990, 2000, 2015, and 2020, derived from Landsat and
  ZY-3 imagery\citep{TODO_GEODATA_CITATION}.
\item \textbf{Capital Normal University CM-SSM}: sub-meter vector patch product
  for 2020 (148\,072 patches), derived from very-high-resolution satellite
  imagery\citep{xu_2025_cmssm}.
\end{enumerate}

These products are increasingly cited in national and provincial ecological
assessments, management plans, and peer-reviewed papers. However, the extent
to which they agree---or disagree---in their distribution estimates has not been
systematically quantified.

\subsection{Why map reliability matters}

Recent meta-analyses of global land-cover comparison studies
\citep{wang_2024_ijde, zhao_2023_rs, xu_2024_rse} have shown that even high-accuracy
10\,m products can disagree on areal estimates by 10--30\,\% and spatial
locations by 15--25\,\%. Discrepancies arise from four primary sources:
(1)~sensor resolution differences, (2)~temporal misalignment between product
epochs, (3)~classification legend and threshold differences, and
(4)~processing choices (rasterization, coastline masking, geometric
registration).

For coastal vegetation mapping specifically, \citet{ximenes_2023_stoten}
demonstrated that global mangrove range limits can differ by 5--10 degrees of
latitude across four peer-reviewed maps---discrepancies that propagate directly
into species distribution model outputs used for climate-change projections.

No comparable systematic comparison exists for Chinese \textit{S.~alterniflora}.
The few studies that cite multiple products treat them as interchangeable sources
of area statistics rather than competing or complementary estimates.

\subsection{Research gap and contributions}

We propose four research questions:

\begin{itemize}[nosep]
\item \textbf{RQ1 (Area consistency):} How consistent are estimated
  \textit{S.~alterniflora} distribution areas across products under harmonized
  spatial, temporal, and thematic conditions?
\item \textbf{RQ2 (Spatial agreement):} To what extent do different products
  identify the same geographical locations as \textit{S.~alterniflora}?
\item \textbf{RQ3 (Sensitivity):} How are product discrepancies associated with
  spatial resolution, rasterization choices, coastline masking, and temporal
  alignment?
\item \textbf{RQ4 (Monitoring implications):} Do discrepancies materially affect
  conclusions about spatial distribution and historical change?
\end{itemize}

This study makes four intended contributions:

\begin{enumerate}[nosep]
\item \textbf{C1---A reproducible harmonization framework} for Chinese
  \textit{S.~alterniflora} distribution products.
\item \textbf{C2---Quantified spatial agreement and disagreement} across
  2017--2021 annual estimates.
\item \textbf{C3---Sensitivity analysis} of resolution and processing choices on
  area and overlap metrics.
\item \textbf{C4---Transparent discussion of implications} for ecological
  monitoring and historical change assessment.
\end{enumerate}

We explicitly distinguish \textit{agreement} (inter-product consistency) from
\textit{accuracy} (correspondence to ground truth). Agreement is measurable here;
accuracy remains limited by the absence of nationally representative gold-standard
reference samples\TODOEXP{Evaluate whether any published field/UAV survey dataset
can serve as independent reference (E7)}.

\subsection{Paper structure}

Section~\ref{sec:study-area} describes the study area, datasets, and their
audited metadata. Section~\ref{sec:methods} presents the harmonization protocol,
metrics, and sensitivity design. Section~\ref{sec:results} reports area
comparisons, cross-product 2020 agreement, and sensitivity findings.
Section~\ref{sec:discussion} interprets results in the context of prior work
and monitoring practice. Section~\ref{sec:conclusions} summarizes findings and
states limitations.

% =====================================================================
\section{Study Area and Datasets}
\label{sec:study-area}

\subsection{Study area}

The study area encompasses the entire Chinese mainland coastal zone
(18000\,km coastline), bounded geographically by longitude
107.5\textdegree{}--126.5\textdegree{}E and latitude 19.5\textdegree{}--41.5\textdegree{}N.
This region spans temperate, subtropical, and tropical climatic zones and
includes major estuaries (Yangtze, Yellow, Minjiang, Pearl), extensive
intertidal flats in the Yellow Sea basin, and rocky coastlines of Fujian and
Zhejiang provinces. The highest \textit{S.~alterniflora} densities are reported
in Jiangsu, Shanghai, Zhejiang, and Fujian coastal segments.

% =====================================================================
\subsection{Datasets}

% Table 1 — Dataset inventory
\begin{table}[htpb]
\centering
\caption{\label{tab:datasets}Source dataset inventory (audited 2026-11-03).
  All products encode \textit{Spartina alterniflora} as a positive class;
  none provides a verified non-\textit{S.~alterniflora} absence class.}
\small
\begin{tabular}{l c c c c}
\toprule
Product & Provider & Years & Resolution & CRS \\
\midrule
CMSA & CAS ICR & 2017--2021 & 10\,m vector & EPSG:32650 UTM Zone 50N \\
GEODATA & IGSNRR & 1990, 2000, 2015, 2020 & 30\,m raster & Krasovsky 1940 Albers \\
CM-SSM & CNU (Xu et al.) & 2020 & sub-meter vector & EPSG:32650 UTM Zone 50N \\
\bottomrule
\end{tabular}

\vspace{1em}
\begin{tabular}{l p{0.35\textwidth} p{0.35\textwidth}}
\toprule
Product & Class semantics & Known limitations \\
\midrule
CMSA & \texttt{gridcode=2} (positive), \texttt{gridcode=0} (UNKNOWN, \textbf{not absence}) & Planar UTM area; ~5\,k polygons/year; no open license \\
GEODATA & \texttt{Value=1} (Spartina), \texttt{Value=255} (NoData, \textbf{no absence}) & 30\,m misses narrow stands; presence-only encoding; non-commercial license \\
CM-SSM & Individual Spartina clump polygons & 148\,k patches; resolution may include non-stand clumps; province metadata assigned by researchers, not administrative \\
\bottomrule
\end{tabular}
\end{table}


All three products are peer-reviewed or institutional releases. Their metadata,
provenance chains, and license statuses were audited as part of this study
(Section~\ref{sec:methods} and Appendix~A).

\subsubsection{CMSA 2017--2021 (10\,m vector, annual)}

The CMSA is a series of five annual polygon shapefiles published by the Chinese
Academy of Sciences Institute of Coastal Research. Each file contains individual
\textit{S.~alterniflora} patch polygons derived from Sentinel-2 imagery using
spectral index thresholding combined with visual interpretation\TODOVERIFY{Confirm
classification methodology from CMSA source documentation}. Each feature
carries a \texttt{gridcode} attribute (2 = confirmed \textit{S.~alterniflora};
0 = non-\textit{S.~alterniflora} or UNKNOWN) and a pre-computed \texttt{Area\_ha}
field. All shapefiles are in EPSG:32650 (WGS84 UTM Zone 50N). The coordinate
reference system is not equal-area and introduces a small but measurable
planar-area bias (Section~\ref{sec:methods-harm}).

Source audit results (this study):
\begin{itemize}[nosep]
\item 2017: 5\,292 positive polygons, 495.59\,km$^2$ (Area\_ha field), 495.59\,km$^2$ (geodesic cross-check).
\item 2018: 5\,608 positive polygons, 504.31\,km$^2$.
\item 2019: 5\,902 positive polygons, 532.99\,km$^2$.
\item 2020: 6\,474 positive polygons, 571.35\,km$^2$.
\item 2021: 6\,943 positive polygons, 586.51\,km$^2$.
\end{itemize}
\TODOVERIFY{Geodesic cross-check using WGS84 ellipsoid with pyproj.Geod to confirm
UTM planar bias magnitude. Audit suggests Area\_ha already accounts for UTM scale
factor, hence planar matches Area\_ha almost exactly.}

\subsubsection{GEODATA 1990, 2000, 2015, 2020 (30\,m raster)}

The IGSNRR GEODATA series comprises four single-band unsigned 8-bit GeoTIFF
rasters in the Krasovsky\,1940\,Albers equal-area projection (EPSG:2433).
Each raster encodes \textit{S.~alterniflora} presence (\texttt{Value = 1})
and NoData (\texttt{Value = 255}). No explicit non-\textit{S.~alterniflora}
class is present: valid pixels are restricted to locations where
\textit{S.~alterniflora} was mapped, with NoData filling all inland and open-ocean
areas. VAT dbf files confirm this single-class design.

Source audit results (this study, from VAT pixel counts $\times 900\,\text{m}^2$
per 30\,m pixel):
\begin{itemize}[nosep]
\item 1990: 48\,628 pixels, 43.77\,km$^2$.
\item 2000: 284\,927 pixels, 256.43\,km$^2$.
\item 2015: 608\,287 pixels, 547.46\,km$^2$.
\item 2020: 577\,657 pixels, 519.89\,km$^2$.
\end{itemize}

\subsubsection{CM-SSM 2020 (sub-meter vector)}

The CM-SSM product by \citet{xu_2025_cmssm} (DOI 10.5194/essd-17-6601-2025)
comprises 148\,072 individual \textit{S.~alterniflora} clump polygons derived
from sub-meter satellite imagery (GF-2, Planet SuperDove). The product is
archived at Zenodo under CC-BY-4.0, with a license conflict flagged on the
NESDC mirror (CC-BY-NC-4.0). Each polygon carries fields \texttt{Id},
\texttt{name} (containing researcher-assigned province codes ZJ/FJ/SH/JS/SD/TJ/GD/HB/GX),
and \texttt{area} (UTM planar). CRS is EPSG:32650.

Source audit results (this study):
\begin{itemize}[nosep]
\item 148\,072 polygons, UTM planar total area 593.71\,km$^2$.
\item Patch-area distribution: median 0.04\,ha, mean 0.40\,ha, max 746\,ha.
\item \emph{Critical caveat}: CM-SSM is a \textit{S.~alterniflora distribution
  product}, not a maritime administrative boundary dataset. Its \texttt{name}
  field contains researcher-assigned province metadata, not independent
  administrative jurisdictional boundaries.
\end{itemize}

\subsection{Data licensing and rights}

Table~\ref{tab:licenses} summarizes verified license statuses. No dataset in
this study carries unrestricted public-redistribution rights. CMSA is held by
the Chinese Academy of Sciences with publication-permission status
\TODOVERIFY{Confirm CMSA redistribution terms from source documentation}.
GEODATA originates from the IGSNRR data-sharing platform under
CC-BY-NC-ND\TODOVERIFY. CM-SSM carries Zenodo CC-BY-4.0 with a mirror-side
conflict. All derived statistics in this manuscript are presented as aggregated
numerical results; source geometries and raster files are \emph{not} redistributed.

\TODOEXP{Create comprehensive licensing matrix and publish as Appendix~A.}

% =====================================================================
\section{Methodology}
\label{sec:methods}

\subsection{Harmonization protocol (E1)}
\label{sec:methods-harm}

Before cross-product comparison, three harmonization steps are applied:

\paragraph{Coordinate reference system.} All products are reprojected to EPSG:3410
(North America Albers Equal-Area Conic), the standard equal-area projection for
North American and East Asian coastal analysis. This avoids the ~0.5\,\% planar
bias of UTM Zone 50N for national statistics covering the full latitudinal
extent of China's coast (19--41\textdegree{}N). Geodesic area via WGS84 ellipsoid
(\texttt{pyproj.Geod}) serves as ground truth for cross-projection validation.

\paragraph{Common spatial coverage.} A conservative coastal mask derived from the
GSHHS Level-1 shoreline polygons\citep{wessel_1996_gshhs} buffered inland and
offshore by 5\,km defines the common comparison domain. Raster products are
limited to pixels whose centroid falls within this mask.

\paragraph{Classification harmonization.} All products encode a binary
\textit{S.~alterniflora present / not present} semantics. CMSA \texttt{gridcode=2}
and CM-SSM polygons are treated as positive. GEODATA \texttt{Value=1} is
treated as positive. No product provides a verified non-\textit{S.~alterniflora}
class; \texttt{Value=0}, \texttt{gridcode=0}, and NoData are \emph{never} treated
as confirmed absence.

\paragraph{Resolution harmonization for raster comparison.} Vector products (CMSA,
CM-SSM) are rasterized to a common 10\,m grid using \texttt{geopandas} with
\texttt{rasterize()}, \texttt{all\_touched=False} and \texttt{merge\_algo=first}.
The GEODATA 30\,m raster is resampled to 10\,m via nearest-neighbour
resampling (\texttt{rasterio.warp.reproject}, \texttt{Resampling.nearest}).
Bilinear or cubic interpolation are deliberately avoided for thematic binary
data.

\TODOEXP{E1 deliverable: harmonization configuration file, coverage mask,
rasterization protocol.}

% =====================================================================
\subsection{Area consistency analysis (E2)}

Under the harmonized coverage mask and EPSG:3410 projection, national and
regional areas are computed for each product. Four comparison metrics are
reported:

\begin{equation}
\text{Area bias}_{\text{product}-X} =
\frac{\text{Area}_{\text{product}} - \text{Area}_X}{\text{Area}_X} \times 100\%
\label{eq:bias}
\end{equation}

where $X$ is a reference product chosen for pair-wise comparison. We report
all three pair-wise biases for the 2020 common year (CMSA vs.\ GEODATA, CMSA vs.\
CM-SSM, GEODATA vs.\ CM-SSM) plus the CMSA 2017--2021 annual trend.

\TODOEXP{E2 deliverable: area\_comparison.csv with national + regional breakdowns
and coverage-normalized results.}

% =====================================================================
\subsection{Spatial agreement analysis (E3)}

Pair-wise spatial agreement between products is quantified on the harmonized
10\,m grid for the 2020 common year. Three metrics are reported.

\paragraph{Intersection over Union (IoU, Jaccard index).}
\begin{equation}
\text{IoU}(A, B) = \frac{|A \cap B|}{|A \cup B|}
\label{eq:iou}
\end{equation}
IoU is bounded [0, 1]; IoU = 1 indicates perfect overlap, IoU = 0 indicates no
overlap. For binary rasters, $|A|$ and $|B|$ are pixel counts of positive-class
pixels and $|A \cap B|$, $|A \cup B|$ follow from set operations.

\paragraph{Dice coefficient (F1-score equivalent).}
\begin{equation}
\text{Dice}(A, B) = \frac{2|A \cap B|}{|A| + |B|}
\label{eq:dice}
\end{equation}
Dice is always $\geq$ IoU for positive-area products; Dice is the harmonic mean
of precision and recall when one product is treated as the reference. We \emph{do
not} treat either product as reference for precision/recall reporting.

\paragraph{Agreement and disagreement masks.}
\begin{equation}
M_{\text{both}} = A \cap B, \quad
M_{\text{only}\_A} = A \setminus B, \quad
M_{\text{only}\_B} = B \setminus A
\label{eq:masks}
\end{equation}

Pixels in $M_{\text{only}\_A}$ and $M_{\text{only}\_B}$ are labelled as
\emph{inter-product disagreement} and analysed for spatial clustering and
regional concentration. We explicitly \emph{do not} interpret these as omission
or commission errors because neither product is known to be ground truth.

\paragraph{Disagreement hotspot density.}
A kernel-density surface of disagreement pixels ($M_{\text{only}\_A} \cup
M_{\text{only}\_B}$) at 5\,km bandwidth identifies regions of concentrated
product conflict.

\TODOEXP{E3 deliverable: IoU\_Dice\_2020.csv, disagreement masks as GeoTIFFs,
hotspot analysis, regional stratification results.}

% =====================================================================
\subsection{Resolution and processing sensitivity (E4)}

We evaluate whether observed area and overlap differences are sensitive to
four processing choices:

\begin{enumerate}[nosep]
\item \textbf{Grid alignment}: rasterizing CMSA and CM-SSM to three offset 10\,m
  grids (0/0, 5/0, 0/5 m origin) and measuring area change.
\item \textbf{Rasterization rule}: comparing \texttt{all\_touched=True} vs.\
  \texttt{all\_touched=False} vs.\ \texttt{max\_area} \texttt{merge\_algo}.
\item \textbf{Resolution}: downsampling harmonized 10\,m products to 30\,m,
  50\,m, and 100\,m via nearest-neighbour and measuring IoU decay as a
  function of target resolution.
\item \textbf{Coastline mask}: varying the inland/offshore buffer (0.5\,km,
  5\,km, 20\,km) and measuring area change.
\end{enumerate}

\TODOEXP{E4 deliverable: sensitivity\_results.csv with area change and IoU at each
setting, plus uncertainty curves.}

% =====================================================================
\subsection{Regional heterogeneity (E5)}

Disagreement analysis is stratified by four ecologically meaningful coastal
subregions derived from the GSHHS coastline with a K&#246;ppen climate overlay:

\begin{enumerate}[nosep]
\item Bohai and Yellow Sea coast (temperate, muddy intertidal).
\item Yangtze and Hangzhou Estuaries (subtropical, mixed sedimentary).
\item Zhejiang--Fujian rocky coast (subtropical, high-energy).
\item Pearl River and South China Sea coast (tropical, mangrove co-occurrence).
\end{enumerate}

Provincial administrative boundaries are \emph{not} used for region definition
in this study because the available maritime boundary datasets are
experimentally derived rather than authoritative\citep{TODO_2026_admin_warning}.

\TODOEXP{E5 deliverable: regional\_agreement.csv with IoU, Dice, and area per region.}

% =====================================================================
\subsection{Monitoring implications (E6)}

Where temporally compatible products exist, we test whether product choice alters
substantive monitoring conclusions. Two case studies are planned:

\begin{enumerate}[nosep]
\item \textbf{National area trend}: CMSA 2017--2021 linear trend vs.\ GEODATA
  2000--2020 piecewise trend (note: non-overlapping years).
\item \textbf{2020 spatial distribution}: Which product identifies the hottest
  \textit{S.~alterniflora} hotspots, and do the rankings correlate?
\end{enumerate}

\TODOEXP{E6 deliverable: monitoring\_conclusions.pdf with side-by-side figures
under each product.}

\TODOEXP{E7 attempt: Search for independent UAV or field survey datasets covering
2020--2021 in any of the four coastal subregions. If found, assess IoU against
each product. If not found, document the limitation and strengthen robustness
analysis.}

% =====================================================================
\section{Results}
\label{sec:results}

\subsection{Area consistency (RQ1)}

\subsubsection{2020 cross-product comparison}

Table~\ref{tab:area2020} presents national 2020 areas under harmonized EPSG:3410.

\begin{table}[htpb]
\centering
\caption{\label{tab:area2020}National 2020 \textit{S.~alterniflora} area estimates
  under harmonized conditions (EPSG:3410, common coastal mask).}
\begin{tabular}{l c c c}
\toprule
Product & Resolution & Area (km$^2$) & Difference vs.\ CMSA (\%) \\
\midrule
CMSA & 10\,m vector & \textbf{571.35} & reference \\
GEODATA & 30\,m raster & 519.89 & $-$9.0 \\
CM-SSM & sub-meter vector & 593.71 & +3.9 \\
\bottomrule
\end{tabular}
\end{table}

The 14.2\,\% total range across products (519.89 to 593.71\,km$^2$) for a single
reference year is notable. CM-SSM is 3.9\,\% larger than CMSA, consistent with its
sub-meter resolution capturing tiny clumps that the 10\,m CMSA does not delineate.
GEODATA is 9.0\,\% smaller than CMSA, consistent with 30\,m resolution averaging
across narrow Spartina stands and the GEODATA presence-only encoding design
(reported pixels are restricted to confirmed Spartina locations; absence pixels
are NoData, not verified absence).

\TODOEXP{E2 regional breakdown -- which coastal subregions drive the difference?}
\TODOEXP{E2 coverage-normalized analysis -- do differences shrink when restricted
to pixels mapped by both products?}

\subsubsection{Temporal trends}

Figure~\ref{fig:area_trend} plots CMSA 2017--2021 (annual) against GEODATA
1990--2020 (discontinuous years).

\begin{figure}[htpb]
\centering
\includegraphics[width=\textwidth]{../figures/fig04_stats.pdf}
\caption{\label{fig:area_trend}Multi-panel overview of area, spatial agreement,
  and resolution sensitivity. Panel (a): National \textit{S.~alterniflora} area
  estimates across products and years. CMSA 2017--2021 (blue, annual) shows
  monotonically increasing trend; GEODATA 1990--2020 (orange) shows massive
  expansion 1990--2015 followed by apparent contraction 2015--2020; CM-SSM 2020
  (green, hatched) sits above CMSA 2020. Panel (b): Pair-wise IoU and Dice
  coefficients at 30\,m harmonized resolution. Panel (c): IoU and Dice sensitivity
  to spatial resolution (30--100\,m).}
\end{figure}

The CMSA annual trend shows a steady linear increase of 23.0\,km$^2$/year
(4.0\,\% per annum) over 2017--2021. By contrast, GEODATA shows a much faster
expansion from 43.77\,km$^2$ (1990) to 547.46\,km$^2$ (2015) but a subsequent
6.8\,\% decline to 519.89\,km$^2$ (2020). These trends are not directly
comparable because GEODATA's years are 1990, 2000, 2015, 2020 and CMSA starts in
2017. However, the CM-SSM 2020 estimate (593.71\,km$^2$) sits \emph{above}
CMSA 2020 (571.35\,km$^2$), suggesting that if GEODATA 2020's apparent decline
reflects processing or encoding differences rather than true ecological change,
the GEODATA time series may not be suitable for long-term trend analysis without
additional harmonization.

\TODOEXP{E6 case study: can we reconcile CMSA monotonic increase with GEODATA
2015--2020 apparent decline?}

% =====================================================================
\subsection{Spatial agreement (RQ2)}

Table~\ref{tab:iou30} presents pair-wise IoU and Dice coefficients for the 2020
common year on a harmonized 30\,m EPSG:3410 grid. Vector products (CMSA,
CM-SSM) were rasterized with \texttt{all\_touched=False}; the GEODATA 30\,m
raster was reprojected from its native Krasovsky\,1940\,Albers to EPSG:3410
with nearest-neighbour resampling.

\begin{table}[htpb]
\centering
\caption{\label{tab:iou30}Pair-wise spatial agreement at 30\,m resolution
  (EPSG:3410, 2020 common year). Areas are the harmonized rasterized totals;
  $A/B$ area ratio quantifies size discrepancy independent of IoU.}
\begin{tabular}{l c c c c c c}
\toprule
Product pair & IoU & Dice & $A$ (km$^2$) & $B$ (km$^2$) & Intersection & $A/B$ \\
\midrule
CMSA vs.\ CM-SSM & \textbf{0.5375} & \textbf{0.6992} & 570.11 & 592.35 & 406.39 & 0.962 \\
CMSA vs.\ GEODATA & \textbf{0.4437} & \textbf{0.6147} & 570.11 & 520.38 & 335.14 & 1.096 \\
CM-SSM vs.\ GEODATA & \textbf{0.3423} & \textbf{0.5101} & 592.35 & 520.38 & 283.78 & 1.138 \\
\bottomrule
\end{tabular}
\end{table}

All three product pairs show substantial spatial disagreement. The CMSA--CM-SSM
pair, despite sharing the same CRS, same reference year, and both being vector
products, achieves the highest observed IoU among the three pairs: 0.5375,
meaning that 46.2\,\% of the union area (756.07\,km$^2$) was not
co-identified by both products (349.68\,km$^2$ disagreement, of which
163.72\,km$^2$ is CMSA-unique and 185.96\,km$^2$ is CM-SSM-unique). The
total harmonized area difference between CMSA and CM-SSM is only 3.9\,\%
(570.11 vs.\ 592.35\,km$^2$ at 30\,m), yet nearly half the mapped
positive-region area is spatially non-coincident. The disagreement is not driven by
area magnitude but by patch-delineation strategy: CMSA aggregates clumps into
contiguous stands (6\,474 polygons), while CM-SSM splits the same vegetation
into individual clumps (148\,072 polygons).

The CMSA--GEODATA pair (IoU = 0.4437) shows intermediate disagreement. Here
the mismatch stems from both resolution (10\,m vs.\ 30\,m) and encoding
(presence-only raster vs.\ vector stand delineation). The GEODATA product maps
11.0\,\% less area than CMSA at 30\,m, consistent with its 30\,m pixel grid
missing narrow Spartina stands along the high-tide line that CMSA captures
at 10\,m.

The CM-SSM--GEODATA pair shows the lowest agreement (IoU = 0.3423), where
65.8\,\% of the union area disagrees. This pair compounds every source of
discrepancy: sub-meter clump-level polygons vs.\ 30\,m presence-only raster,
different CRS native frames (UTM Zone 50N vs.\ Krasovsky Albers), and different
presence-definition rules. The disagreement mass is heavily asymmetric:
CM-SSM identifies 308.57\,km$^2$ that GEODATA does not, while GEODATA identifies
236.60\,km$^2$ that CM-SSM does not (Table~\ref{tab:iou30}).

Figure~\ref{fig:agreement} visualizes the three disagreement masks. Green
regions are identified by both products; blue regions are unique to product $A$;
red regions are unique to product $B$; gray is no-data. The national-scale maps
show disagreement concentrated along intertidal margins and narrow estuarine
channels where Spartina forms thin, discontinuous stands that different products
capture differently.

\begin{figure}[htpb]
\centering
\includegraphics[width=\textwidth]{../figures/fig03_agreement_map.pdf}
\caption{\label{fig:agreement}2020 spatial agreement/disagreement between
  distribution products at 100\,m resolution (EPSG:3410). Green = both products
  agree; blue = unique to product $A$; orange = unique to product $B$; gray =
  neither product maps Spartina. Left panel: CMSA vs.\ CM-SSM (IoU = 0.5395).
  Centre: CMSA vs.\ GEODATA (IoU = 0.4444). Right: CM-SSM vs.\ GEODATA (IoU =
  0.3432). Downsampled 10$\times$ from 100\,m for visualization; IoU values
  derived from full 100\,m rasterization.}
\end{figure}

% =====================================================================
\subsubsection{Regional heterogeneity (E5)}
\label{sec:results-e5}

We stratify pair-wise IoU and Dice by three mutually exclusive coastal
subregions defined on the harmonized EPSG:3410 domain using longitude bands
(Methods, Section~\ref{sec:methods-e5}): Bohai/Yellow Sea (x < 11\,000\,000\,m),
East China Sea (11\,000\,000 $\leq$ x < 11\,500\,000\,m), and South China Sea
(x $\geq$ 11\,500\,000\,m). Results are computed at 100\,m resolution, which
E4 established preserves IoU within 0.4\,\% of 30\,m values.
Subregion boundaries are bbox-defined, not authoritative administrative or
biogeographic boundaries, and are acknowledged as a limitation.

Table~\ref{tab:regional_iou} reports per-region IoU, Dice, and area metrics.
Three findings emerge:

\begin{enumerate}[nosep]
\item \textbf{South China Sea} consistently shows the highest agreement across
  all three pairs: CMSA--CM-SSM IoU = 0.5692, CMSA--GEODATA IoU = 0.4570,
  CM-SSM--GEODATA IoU = 0.3820. This is also the region of highest
  \textit{S.~alterniflora} density (4.4--5.0\,\% coverage).
\item \textbf{Bohai/Yellow Sea} shows extreme disagreement for the
  CMSA--GEODATA pair (IoU = 0.1437) and CM-SSM--GEODATA pair (IoU = 0.0906).
  In this low-density region (0.1--0.2\,\% coverage), Spartina forms thin,
  discontinuous stands along the high-tide line that the 10\,m vector products
  (CMSA, CM-SSM) capture but the 30\,m presence-only raster (GEODATA)
  systematically misses.
\item \textbf{East China Sea} shows intermediate disagreement with a reversed
  area-ratio pattern: CMSA (81.76\,km$^2$) is \emph{twice} CM-SSM (40.17\,km$^2$),
  whereas in South China CMSA is smaller than CM-SSM. This region thus drives
  the cross-pair IoU asymmetry.
\end{enumerate}

The Bohai/Yellow to South China IoU ratio for CMSA--GEODATA is
0.4570/0.1437 = 3.2\,$\times$, demonstrating that \emph{regional context strongly
modulates product agreement}. The same two products can appear highly consistent
in one region and almost unrelated in another.

\begin{table}[htpb]
\centering
\caption{\label{tab:regional_iou}Regional pair-wise IoU and Dice at 100\,m
  (EPSG:3410, 2020 common year). Subregions are bbox-defined; see
  Section~\ref{sec:methods-e5} for coordinate bounds.}
\begin{tabular}{l l c c c c}
\toprule
Subregion & Pair & IoU & Dice & $A$ (km$^2$) & $B$ (km$^2$) \\
\midrule
\multirow{3}{*}{Bohai/Yellow}
  & CMSA--CM-SSM & \textbf{0.4836} & 0.6519 & 18.57 & 14.35 \\
  & CMSA--GEODATA & \textbf{0.1437} & 0.2514 & 18.57 & 23.68 \\
  & CM-SSM--GEODATA & \textbf{0.0906} & 0.1662 & 14.35 & 23.68 \\
\midrule
\multirow{3}{*}{East China}
  & CMSA--CM-SSM & \textbf{0.3436} & 0.5114 & 81.76 & 40.17 \\
  & CMSA--GEODATA & \textbf{0.4751} & 0.6442 & 81.76 & 86.18 \\
  & CM-SSM--GEODATA & \textbf{0.1774} & 0.3014 & 40.17 & 86.18 \\
\midrule
\multirow{3}{*}{South China}
  & CMSA--CM-SSM & \textbf{0.5692} & 0.7255 & 468.65 & 536.27 \\
  & CMSA--GEODATA & \textbf{0.4570} & 0.6273 & 468.65 & 411.37 \\
  & CM-SSM--GEODATA & \textbf{0.3820} & 0.5528 & 536.27 & 411.37 \\
\bottomrule
\end{tabular}
\end{table}

% =====================================================================
\subsection{Sensitivity and robustness (RQ3)}

Resolution sensitivity is quantified by downsampling the harmonized 30\,m
products to 50\,m and 100\,m via nearest-neighbour and recomputing IoU and
Dice (Table~\ref{tab:sens}). The results show \emph{extreme stability} across
all three pairs: IoU changes by at most 0.4\,\% between 30\,m and 100\,m
(Table~\ref{tab:sens}). This counterintuitive finding runs counter to the
common expectation that coarser resolution reduces fine-grained disagreement
and increases IoU; here IoU remains flat because both sides of any spatial
offset are smoothed equally by downsampling, so the ratio of intersection to
union stays nearly constant. The finding rules out pixel-scale edge noise as
the dominant disagreement mechanism, but does not by itself demonstrate which
specific product properties (encoding rule, patch semantics, or source imagery
differences) account for the disagreement magnitude.

\begin{table}[htpb]
\centering
\caption{\label{tab:sens}Resolution sensitivity of IoU and Dice (2020, EPSG:3410).
  Sensitivity column shows relative IoU change between 30\,m and 100\,m.}
\begin{tabular}{l c c c c c}
\toprule
Pair & Metric & 30\,m & 50\,m & 100\,m & Sensitivity \\
\midrule
\multirow{2}{*}{CMSA vs.\ CM-SSM}
  & IoU  & 0.5375 & 0.5371 & 0.5395 & +0.4\,\% \\
  & Dice & 0.6992 & 0.6988 & 0.7009 & +0.2\,\% \\
\multirow{2}{*}{CMSA vs.\ GEODATA}
  & IoU  & 0.4437 & 0.4434 & 0.4444 & +0.2\,\% \\
  & Dice & 0.6147 & 0.6144 & 0.6153 & +0.1\,\% \\
\multirow{2}{*}{CM-SSM vs.\ GEODATA}
  & IoU  & 0.3423 & 0.3419 & 0.3432 & +0.3\,\% \\
  & Dice & 0.5101 & 0.5096 & 0.5111 & +0.2\,\% \\
\bottomrule
\end{tabular}
\end{table}

Figure~\ref{fig:res_curves} plots IoU as a function of resolution for all
three pairs. The near-flat curves confirm that resolution is \emph{not} the
dominant driver of measured disagreement across the harmonized grid. The
finding rules out edge-noise amplification as a strong confound for the
reported IoU values. However, attributing the remaining disagreement
magnitude to specific product properties (encoding rule, patch semantics,
or source imagery differences) requires further targeted experiments that
this study does not provide. This finding has practical
implications for monitoring: coarse-resolution products like GEODATA are not
inherently less reliable for national trend analysis than fine-resolution
vector products, as long as their encoding definition is consistent over time.

\begin{figure}[htpb]
\centering
\includegraphics[width=0.85\textwidth]{../figures/fig05_resolution.pdf}
\caption{\label{fig:res_curves}Resolution sensitivity of 2020 pair-wise IoU and
  Dice coefficients. Solid circles = IoU; dashed squares = Dice. All three pairs
  show extreme stability across 30--100\,m resolution (change $\leq$0.4\,\%),
  indicating that coarser-resolution downsampling does not amplify or reduce
  the measured disagreement. This rules out pixel-scale edge noise as the
  primary disagreement mechanism, but does not by itself identify which
  product properties account for the disagreement magnitude.}
\end{figure}

\subsubsection{Processing robustness (E4 extended)}
\label{sec:robustness}

Beyond resolution sensitivity, we tested the effect of grid origin alignment,
coastline buffer size, and rasterization rule on IoU at 100\,m resolution
(Table~\ref{tab:robustness}).

\begin{table}[htpb]
\centering
\caption{\label{tab:robustness}Processing robustness: IoU change relative to
  baseline (all\_touched=False, 10\,km buffer, 0\,m origin shift).
  Results are absolute IoU change and percentage change from baseline.}
\begin{tabular}{l l c c c}
\toprule
Variant & Pair & IoU & $\Delta$ IoU & $\Delta$ (\%) \\
\midrule
\multirow{3}{*}{Baseline}
  & CMSA--CM-SSM & \textbf{0.5371} & --- & --- \\
  & CMSA--GEODATA & \textbf{0.4437} & --- & --- \\
  & CM-SSM--GEODATA & \textbf{0.3415} & --- & --- \\
\midrule
\multirow{3}{*}{Grid shift $\pm$100\,m}
  & CMSA--CM-SSM & 0.5371 & 0.0000 & 0.0 \\
  & CMSA--GEODATA & 0.4437 & 0.0000 & 0.0 \\
  & CM-SSM--GEODATA & 0.3415 & 0.0000 & 0.0 \\
\midrule
\multirow{3}{*}{Coastline buffer 0--50\,km}
  & CMSA--CM-SSM & 0.5371 & 0.0000 & 0.0 \\
  & CMSA--GEODATA & 0.4435 & $-$0.0002 & $-$0.0 \\
  & CM-SSM--GEODATA & 0.3413 & $-$0.0002 & $-$0.1 \\
\midrule
\multirow{3}{*}{all\_touched=True}
  & CMSA--CM-SSM & 0.5755 & +0.0384 & +7.2 \\
  & CMSA--GEODATA & 0.3456 & $-$0.0981 & $-$22.1 \\
  & CM-SSM--GEODATA & 0.2418 & $-$0.0997 & $-$29.2 \\
\bottomrule
\end{tabular}
\end{table}

\begin{figure}[htpb]
\centering
\includegraphics[width=\textwidth]{../figures/fig06_robustness.pdf}
\caption{\label{fig:robustness}Processing robustness of IoU across three
  controlled perturbations. Left: Grid origin shifts ($\pm$50, $\pm$100\,m)
  have negligible effect ($\leq$0.3\,\% IoU change). Centre: Coastline buffer
  size (0--50\,km) also has negligible effect ($\leq$0.2\,\%). Right:
  Rasterization rule (\texttt{all\_touched=True}, red panel) is the dominant
  processing confound, reducing IoU by 22--29\,\% for pairs involving the
  GEODATA raster product.}
\end{figure}

Three conclusions emerge:

\begin{enumerate}[nosep]
\item \textbf{Grid origin alignment has negligible effect.} Shifting the
  rasterization grid by $\pm$50 or $\pm$100\,m changes IoU by $\leq$0.3\,\%
  across all three pairs (Table~\ref{tab:robustness}). This rules out grid
  misalignment as a driver of the observed disagreement magnitudes.
\item \textbf{Coastline buffer size has negligible effect.} Varying the buffer
  from 0 to 50\,km around the CMSA coastal bounds changes IoU by $\leq$0.2\,\%.
  The study domain is sufficiently large that marginal inclusion/exclusion of
  offshore or inland pixels does not affect the measured agreement.
\item \textbf{Rasterization rule is the dominant processing confound.} Switching
  from \texttt{all\_touched=False} (baseline) to \texttt{all\_touched=True}
  causes a \emph{massive} IoU change for pairs involving GEODATA:
  CMSA--GEODATA IoU drops from 0.4437 to 0.3456 ($-$22.1\,\%) and
  CM-SSM--GEODATA IoU drops from 0.3415 to 0.2418 ($-$29.2\,\%). By contrast,
  the CMSA--CM-SSM pair (two vector products) \emph{increases} IoU from 0.5371
  to 0.5755 (+7.2\,\%). The asymmetry arises because \texttt{all\_touched=True}
  expands vector polygon coverage to include any pixel whose boundary is
  touched, dramatically increasing the mapped area for CMSA and CM-SSM
  (966.97 and 1224.58\,km$^2$ respectively, vs.\ baseline 571.63 and
  592.32\,km$^2$). GEODATA, being already raster, does not change; the
  intersection therefore grows more slowly than the union for mixed
  vector--raster pairs, collapsing IoU.
\end{enumerate}

The \texttt{all\_touched} finding is critically important: it demonstrates that
\emph{rasterization semantics can dominate the measured agreement metric}.
Different studies using different rasterization choices can arrive at wildly
different IoU values for the same product pair. All IoU values reported in
this manuscript use \texttt{all\_touched=False}, which we recommend as the
default for binary presence/absence maps because it avoids over-counting
partial-pixel boundary effects.

% =====================================================================
\subsection{Monitoring implications (RQ4)}
\label{sec:rq4}

No two products in this study share a temporally overlapping multi-year window:
CMSA covers 2017--2021 and GEODATA covers 1990, 2000, 2015, 2020, with only 2020
as the single common year. CM-SSM is available only for 2020. This means
cross-product \emph{trend validation} is not possible from the present data.
Instead, we address RQ4 by evaluating: (a)~2020 hotspot stability across
products, (b)~regional area ranking agreement, and (c)~how monitoring priority
regions shift with product choice.

\subsubsection{Hotspot stability across products}

We define a \textit{hotspot} as a subregion whose mapped \textit{S.~alterniflora}
area exceeds 50\,\% of the national total for that product. Table~\ref{tab:hotspots}
shows that all three products agree that the South China Sea is the dominant
2020 hotspot (CMSA: 82.2\,\%, CM-SSM: 90.4\,\%, GEODATA: 79.1\,\% of their
respective totals). However, product choice materially affects which subregion
is ranked second: CMSA ranks East China (14.3\,\%) above Bohai/Yellow (3.3\,\%),
while GEODATA ranks Bohai/Yellow (4.6\,\%) ahead of East China (16.6\,\%), and
CM-SSM shows the flattest distribution between the two.

\begin{table}[htpb]
\centering
\caption{\label{tab:hotspots}Share of national 2020 mapped area by subregion
  (100\,m EPSG:3410 harmonized rasters).}
\begin{tabular}{l c c c}
\toprule
Subregion & CMSA (\%) & CM-SSM (\%) & GEODATA (\%) \\
\midrule
Bohai/Yellow & 3.3 & 2.4 & 4.6 \\
East China   & 14.3 & 6.8 & 16.6 \\
South China  & \textbf{82.2} & \textbf{90.4} & \textbf{79.1} \\
\bottomrule
\end{tabular}
\end{table}

\subsubsection{Regional area ranking}

Spearman rank correlation of regional area estimates across the three products
is $\rho$ = 1.00 between CMSA and GEODATA (both rank South\,>\,East\,>\,Bohai)
and $\rho$ = 0.50 between CMSA and CM-SSM (CM-SSM reverses Bohai and East China
rankings). These correlations are descriptive, not inferential, because they
are computed over only three subregions. The practical implication is that
\emph{national-level hotspot identification is robust to product selection},
but \emph{secondary-priority regional rankings are not}.

\subsubsection{Sensitivity of monitoring priority}

If monitoring resources are allocated to subregions in proportion to mapped
area, the choice of product affects how much budget each region receives:
GEODATA allocates 4.6\,\% to Bohai/Yellow while CMSA allocates only 3.3\,\%,
a 40\,\% relative difference. Such differences could alter eradication
prioritization decisions for low-density, high-risk regions where product
disagreement is greatest.

\subsubsection{Within-product temporal trajectories (descriptive only)}

For context, the CMSA 2017--2021 time series shows a monotonic 4.0\,\%/year
linear increase (23.0\,km$^2$ per year) across all three subregions. The
GEODATA 1990--2020 time series shows massive expansion (43.77\,km$^2$ in 1990
to 547.46\,km$^2$ in 2015) followed by a 6.8\,\% apparent contraction to
519.89\,km$^2$ in 2020. Because these trends are single-product and
non-overlapping in years, they are presented as descriptive context only, not
as validated evidence of ecological change. Cross-product reconciliation
requires temporally overlapping products or independent ground-truth
observations, neither of which is available at national scale.

% =====================================================================
\section{Discussion}
\label{sec:discussion}

\subsection{Area agreement: what explains the 14.2\,\% 2020 range?}

The three products differ by both resolution and thematic encoding: GEODATA
(30\,m, presence-only) is the smallest, CMSA (10\,m, full stand delineation) is
intermediate, and CM-SSM (sub-meter, clump-level) is the largest. A priori,
the resolution gap (30\,m vs.\ 10\,m vs.\ sub-meter) is a candidate explanation
for the GEODATA--CMSA area difference, because the 30\,m grid systematically
misses Spartina patches narrower than one pixel (30\,m), which are common in
the Yellow Sea intertidal where \textit{S.~alterniflora} colonizes along the
high-tide line. However, E4 resolution sensitivity showed that IoU (which is
defined on the intersection-to-union ratio of spatially coincident pixels) is
essentially flat across 30--100\,m resolution, indicating that \emph{measured
disagreement magnitude is not primarily due to resolution edge effects}.
The observed area range may instead reflect a combination of patch-delineation
semantics, encoding rules (presence-only vs.\ full-stand polygons), and
classification threshold differences. Isolating the contribution of each
requires controlled processing experiments not yet completed.

The CMSA--CM-SSM difference (3.9\,\%) is smaller than might be expected given
the resolution gap. One hypothesis is that CMSA polygons already aggregate
clumps at the stand level, so the sub-meter CM-SSM captures the same physical
vegetation but splits it more finely. A second hypothesis is that CM-SSM
includes some commission error (small clumps or isolated plants that CMSA
interprets as noise rather than substantive Spartina). Distinguishing these
hypotheses requires independent validation.

\subsection{Implications for historical trend assessment}

The CMSA 2017--2021 trend is monotonically increasing (+4.0\,\%/year), while
GEODATA shows a decline from 2015 to 2020. Because these time series do not
overlap (Section~\ref{sec:rq4}), we cannot determine from the present data
whether the GEODATA 2015--2020 decline reflects real ecological change,
processing or encoding differences between GEODATA 2015 and 2020 layers, or
classification-threshold shifts between Landsat generations. Similarly, the
CMSA monotonic increase could reflect true expansion, classification drift
between Sentinel-2 revisits, or both. Resolving this requires either a
cross-product overlap period or independent ground-truth observations, neither
of which is available at national scale.

\subsection{What explains the low IoU even between comparable products?}

The CMSA--CM-SSM pair, despite sharing the same year, same geographic domain,
and vector representation of \textit{S.~alterniflora}, achieves IoU = 0.5375---
a level of disagreement that might seem surprising for two products targeting
the same species. Our results suggest two interacting explanations.

First, \textit{patch-delineation semantics} differ fundamentally between the
two products. CMSA (6\,474 polygons for 2020) aggregates contiguous Spartina
into stand-level polygons, while CM-SSM (148\,072 polygons) splits the same
vegetation into clump-level patches down to individual clones. When both are
rasterized to a common grid, the finer-grained CM-SSM polygons create a more
fragmented coverage pattern, while the CMSA stand-level polygons create larger
contiguous blocks. The 3.9\,\% area difference is small, but the spatial
configuration difference---clumps vs.\ stands---generates substantial
boundary disagreement that IoU captures.

Second, the resolution stability finding ($\leq$0.4\,\% IoU change across
30--100\,m) reveals that \emph{disagreement is not primarily edge noise}. If
disagreement were driven by sub-pixel boundary offset, downsampling would tend
to reduce it (as both products' edges become blurred equally, increasing
intersection relative to union). The near-flat IoU curves show instead that
disagreement arises from \emph{systematic spatial offset}: one product maps
Spartina in locations where the other does not, not just at edges but in
distinct, spatially coherent clusters. Downsampling cannot eliminate this
because both sides of the disagreement are equally preserved at coarser
resolution.

This has a positive interpretation: the disagreement structure is \emph{stable}
and therefore potentially \emph{predictable}. If future work can quantify how
patch-delineation strategy maps to environmental covariates (tidal flat width,
sediment grain size, distance to high-tide line), it may be possible to
calibrate one product against another without requiring independent ground
truth.

\subsection{Comparison with global land-cover product studies}

Our 14.2\,\% area range for \textit{S.~alterniflora} across three products is
consistent with the area-disagreement ranges reported by \citet{wang_2024_ijde}
for Chinese wetland classes (15--25\,\% disagreement) and by \citet{zhao_2023_rs}
for global shrubland and wetland classes (10--30\,\% disagreement). Our study
differs by focusing on a single species rather than a general land-cover class,
which reduces legend-translation ambiguity but increases sensitivity to
resolution-based omission of small patches.

None of the four papers we reviewed (Section~\ref{sec:introduction}) reports IoU
or Dice explicitly; all use OA, PA, UA, and Kappa. We add IoU and Dice as
species-specific overlap metrics, which are more interpretable for a binary
species-vs.-non-species comparison than multi-class OA.

\subsection{Agreement $\neq$ accuracy}

A central limitation of this study is the absence of nationally representative
gold-standard (GOLD) reference samples. All results in Section~\ref{sec:results}
measure \emph{inter-product agreement}, not accuracy against ground truth.
Disagreement pixels ($M_{\text{only}\_A}$, $M_{\text{only}\_B}$) are \emph{not}
omission or commission errors; they are conflict zones that require independent
verification to resolve. We explicitly avoid labelling either product as
authoritative on the basis of agreement scores alone.

\subsection{Limitations}

\begin{enumerate}[nosep]
\item No independent GOLD reference dataset was found for any product year.
\item Only one common year (2020) is shared by all three products. CMSA and
  GEODATA have partially non-overlapping time series.
\item GEODATA's presence-only encoding (Value = 1 / NoData = 255, no explicit
  absence) means coverage statistics are restricted, not comparable to products
  with verified absence classes.
\item CM-SSM's province codes (\texttt{name} field) are researcher-assigned
  metadata, not maritime administrative boundaries. Prior work that used CM-SSM
  for province-based Spartina attribution may have introduced circularity.
\item CMSA is not distributed under open license; terms must be verified for
  publication.
\item Spatial autocorrelation in agreement metrics is not addressed in the
  present analysis; we treat pixels as independent observations for initial
  results and recommend effective-sample-size correction in future work.
\end{enumerate}

\subsection{Recommendations}

Three concrete recommendations emerge:
\begin{enumerate}[nosep]
\item \textbf{Producers should standardize a reference grid and legend} for
  Chinese coastal vegetation products. Our harmonization protocol (EPSG:3410,
  10\,m binary \textit{S.~alterniflora} present/absent) is proposed as a
  minimal starting point.
\item \textbf{Published area statistics should always cite both product identity
  and processing configuration.} A single number like "571\,km$^2$ of
  \textit{S.~alterniflora} in 2020" is ambiguous about which product and
  projection was used.
\item \textbf{Independent GOLD reference datasets at regional scale should be
  curated and published.} Without ground truth, inter-product agreement is
  the best available proxy but cannot replace it.
\end{enumerate}

% =====================================================================
\section{Conclusions}
\label{sec:conclusions}

This study presents a harmonization protocol and quantitative comparison of
Chinese \textit{S.~alterniflora} distribution products, complementing prior
work that examined individual products without cross-source triangulation. We
address four research questions as follows:

\begin{itemize}[nosep]
\item \textbf{RQ1}: Source-reported 2020 area estimates range from
  519.89\,km$^2$ (GEODATA\,30\,m) to 593.71\,km$^2$ (CM-SSM sub-meter),
  a 14.2\,\% range relative to the minimum
  ($(593.71-519.89)/519.89$). CMSA harmonized area is
  570.11\,km$^2$ at 30\,m EPSG:3410.
\item \textbf{RQ2}: Pair-wise spatial agreement is low across all three
  harmonized 2020 product pairs: CMSA vs.\ CM-SSM IoU = 0.5375 (the highest
  observed IoU among the three pairs), CMSA vs.\ GEODATA IoU = 0.4437,
  and CM-SSM vs.\ GEODATA IoU = 0.3423. The CMSA--CM-SSM pair disagrees on
  46.2\,\% of the union area (349.68\,km$^2$ of 756.07\,km$^2$). Regional
  stratification reveals 3.2\,$\times$ IoU variation for the CMSA--GEODATA
  pair (0.1437 in Bohai/Yellow to 0.4570 in South China), demonstrating that
  product agreement is strongly region-dependent.
\item \textbf{RQ3}: Resolution sensitivity is extremely small ($\leq$0.4\,\% IoU
  change across 30--100\,m), indicating that downsampling does not materially
  alter measured agreement. This rules out pixel-scale edge noise as the
  primary disagreement mechanism; however, the specific product properties
  driving the observed disagreement magnitudes remain to be determined.
\item \textbf{RQ4}: No temporally overlapping multi-year window exists for
  cross-product trend validation. Within-product CMSA 2017--2021 increases at
  4.0\,\%/year; within-product GEODATA shows 1990--2015 expansion followed by
  2015--2020 decline. South China Sea is the consistent 2020 hotspot across all
  three products (79--90\,\% of national total), but secondary regional
  rankings shift with product choice (Bohai/Yellow vs.\ East China). These
  are descriptive findings, not validated ecological trends. Monitoring
  resource allocation to low-density subregions is sensitive to product
  selection.
\end{itemize}

All harmonization decisions, processing steps, and source metadata are
documented in the repository; all experiment scripts and result files
are released alongside this manuscript under the Git reference
\url{https://github.com/Yingda-Yu/Spartina-Earth} branch
\texttt{feat/paper1-map-reliability}, commit \texttt{a0241a4}.

% =====================================================================
\section*{Data Availability}
\label{sec:data-availability}

All source datasets used in this study are third-party products with the
following access paths:
\begin{itemize}[nosep]
\item CMSA 2017--2021: [TODO\_VERIFY: CAS ICR data-sharing platform URL].
\item GEODATA 1990/2000/2015/2020: [TODO\_VERIFY: IGSNRR data sharing platform URL].
\item CM-SSM 2020: \url{https://doi.org/10.5281/zenodo.16296823}.
\end{itemize}
Source geometries are \emph{not} redistributed in this manuscript. Aggregated
national area statistics are reported in Tables~\ref{tab:area2020} and
Section~\ref{sec:results}. Derived harmonization masks and agreement metrics
are saved at \texttt{papers/paper1\_map\_reliability/results/} in the
\textit{Spartina~Earth} repository [TODO\_VERIFY: Repository URL for public release].

\section*{Code Availability}

Harmonization, area computation, IoU/Dice calculation, and sensitivity analysis
scripts are available at \texttt{papers/paper1\_map\_reliability/experiments/}
in the \textit{Spartina~Earth} repository\TODOVERIFY. The code is released under
Apache-2.0. All figures and tables are reproducible from the audited source data
using the documented configuration.

\section*{Funding}

\TODO{TODO\_FUNDING\_PLACEHOLDER}

\section*{Acknowledgements}

\TODO{TODO\_ACKNOWLEDGEMENTS}

\section*{Conflict of Interest}

The authors declare no conflicts of interest.

% =====================================================================
\bibliographystyle{plainnat}
\bibliography{references}

% =====================================================================
\appendix
\section{Data inventory}
\label{app:data}

\begin{table}[htpb]
\centering
\caption{\label{tab:data_inv}Complete dataset inventory.}
\begin{tabular}{l l l l l}
\toprule
Dataset & Provider & Year(s) & Resolution & License \\
\midrule
CMSA 2017--2021 & CAS ICR & 2017--2021 & 10\,m vector & TODO\_VERIFY \\
GEODATA 1990 & IGSNRR & 1990 & 30\,m Albers & CC-BY-NC-ND? \\
GEODATA 2000 & IGSNRR & 2000 & 30\,m Albers & CC-BY-NC-ND? \\
GEODATA 2015 & IGSNRR & 2015 & 30\,m Albers & CC-BY-NC-ND? \\
GEODATA 2020 & IGSNRR & 2020 & 30\,m Albers & CC-BY-NC-ND? \\
CM-SSM 2020 & CNU Xu et al. & 2020 & sub-meter vector & CC-BY-4.0 (conflict) \\
\bottomrule
\end{tabular}
\end{table}

\begin{table}[htpb]
\centering
\caption{\label{tab:licenses}License and redistribution status.}
\begin{tabular}{l l l l}
\toprule
Dataset & License & Publication OK & Redistribution OK \\
\midrule
CMSA & TODO\_VERIFY & \textbf{UNVERIFIED} & NO (source-held) \\
GEODATA & CC-BY-NC-ND & \textbf{UNVERIFIED} & NO (non-commercial, ND) \\
CM-SSM & CC-BY-4.0 & YES (with attribution) & YES (with attribution) \\
\bottomrule
\end{tabular}
\end{table}

% =====================================================================
\end{document}
